%%%%%%%%%%%%%%%%%%%%%%%%%%%%%%%%%%%%%%%%%%%%%%%%%%%%%%%%%%%%%%%%%%%
\section{Workload 1: Deception Detection}
\label{sec:Results}
%%%%%%%%%%%%%%%%%%%%%%%%%%%%%%%%%%%%%%%%%%%%%%%%%%%%%%%%%%%%%%%%%%%
% Restructured Aug 31 per MN's markup: this section now holds ALL the
% deception-detection evaluation -- the experimental setup and audit
% moved here from Sec. IV, along with Table II (tab:audit) and the
% first pipeline's four-domain results.

We first evaluate the mappings of Sec.~\ref{sec:Mapping} on a
deception detection workload, where the
library is set by how many behaviors and domains someone decides to
watch\inote[aiviolet]{Sep 2: your inserts (``a,'' ``workload'') and
the cut of ``on the first of the paper's two workloads:'' applied.
Applying your strikes on ``This section''/``evaluates'' literally
leaves no subject or verb, so the opener is ``We first evaluate,''
taking the ``We 1st'' you wrote above the heading as the cue --
rewrite if you meant something else} -- here, the $K\approx10^{2}$--$10^{3}$ regime of
Fig.~\ref{fig:ladder}. We describe the experimental setup, then
report {\bf (1)} the ten-domain fidelity of each encoding, {\bf (2)}
whether wider codes improve cross-domain transfer, {\bf (3)} how the
two-stage shortlist behaves across bank sizes, and {\bf (4)} the
destroyed-vector audit that disciplines every configuration reported.
Sec.~\ref{sec:LargeK} evaluates the other end of the $K$ axis, where
the library is set by the vocabulary.\inote[aiviolet]{Sep 2: the
``One may reasonably ask why we evaluate at small $K$'' paragraph
(added last round) is deleted per your full strike. Your pink box
next to it -- whether Truthfulness connects to VLA, or to fall back
on the established-dataset argument, or drop the question -- is
parked for a conversation rather than answered in the text; happy to
brainstorm it in chat whenever}

\begin{figure*}[t]
\centering
\includegraphics[width=7.0in]{fig_grids_v2.pdf}
\caption{Cross-domain \auroc, Llama-3.3-70B-Instruct activations at
layer 33. Rows index the domain whose examples are stored in the bank,
columns the domain the held-out test activations come from. Top row,
four-domain published protocol: (a) the full-precision Truthfulness
Spectrum baseline~\cite{ying2026spectrum} as reproduced in our
evaluation; (b) the first-pipeline Euclidean configuration -- 128-d
learned projection, 3-bit LSQ rows, Euclidean matchline
(Sec.~\ref{subsec:FirstPipeline}) -- scored only on datasets shared
with (a); (c) the cell-wise change, (b) minus (a). Bottom row,
ten-domain protocol: (d) 384-bit \signrp ternary code, requiring a
wildcard-capable \tcam cell; (e) 384-bit thermometer code, requiring
only a plain binary \cam cell; (f) Hamming shortlist over the whole
bank followed by exact cosine rerank on raw activations, i.e., a
full-precision ceiling for the same data and protocol. The two
protocols use different datasets and splits and are not comparable
cell by cell.}
\label{fig:grids}
\end{figure*}

\subsection{Experimental setup}
\label{subsec:Setup}
Our evaluation case follows the cross-domain protocol of the
Truthfulness Spectrum study~\cite{ying2026spectrum}: a model produces
honest and deceptive statements across several subject-matter
domains, its activations are recorded, and a detector built from one
domain's examples is scored on every other domain's -- the ``stress
test'' of Sec.~\ref{subsec:Probes}. Activations are taken from
Llama-3.3-70B-Instruct~\cite{llama3herd}. The evaluation grid has ten
domains:
definitional, empirical, fictional, logical, ethical, sycophancy,
\texttt{repe\_honesty}, roleplaying, insider trading, and
sandbagging.\ib{Item 3: our evidential is their empirical and our
ethics\_\_commonsense is their ethical (same source datasets), so the
domains are relabeled to the source paper's names here and in
Fig.~2(d)--(f) (``evid'' $\to$ ``empir''); repe\_honesty is not one of
theirs and stays. Panels (d)--(f) checked cell by cell against the
result files: all three are the current 384-b runs (seed 0, 70/30
split).} The
first four plus sycophancy correspond to the published
grid~\cite{ying2026spectrum}; \texttt{repe\_honesty} follows the
representation-engineering dataset~\cite{zou2023repe}. Within each
domain, 70\% of examples form the stored bank and 30\% are held out for
test. Every cell of the grid is a cross-domain measurement: rows index
the domain whose examples are in the bank, columns the domain the test
activations come from. Bank sizes differ substantially by domain, from
831 rows (definitional) to 3{,}201 (sycophancy), which turns out to
matter in Sec.~\ref{subsec:Shortlist}. Activations are read at layer
33: the source study selected that layer (with logistic-regression
probes on average-token activations) from a cross-domain tuning sweep
on its truth-type datasets, and we inherit the choice so our grids
can be compared against theirs -- i.e., it is a tuned choice for this
model family and scale rather than a transferable constant. All
numbers reported below are from cluster runs on this
configuration.\inote[aiviolet]{Sep 1, per your ``more tutorial in
nature'' note (and the page-4 ``have we even mentioned layer 33''
note, resolved as one fix): the subsection now opens by introducing
the evaluation case, and layer 33 arrives later in the paragraph with
the why -- inherited from the source study's tuning, per their
Sec.~3. Sec.~IV-B now points here instead of introducing the layer
itself}
\ib{Item 4: the dead zone was on. The grid script only has a ternary
encoder (no binary option) and maps the smallest 20\% of projection
magnitudes in each column to don't-care by default, so the 384-b
ternary grid is genuinely ternary and the TCAM-cell claims in Secs.~IV
and V and Table~II stand. One detail still to fix: the don't-care
threshold for queries is computed from the test batch rather than
fixed from training; I will switch it to a training-set threshold and
rerun to confirm the numbers hold.}

\subsection{The grids}
\label{subsec:Grids}
\inote[aiviolet]{Sep 1: the ``measurements here are representative of
the $K\approx10^{2}$--$10^{3}$ rungs'' recap that opened this
subsection is cut per your ``Redundant? Can probably cut'' -- the
reworked section intro above carries it}

Fig.~\ref{fig:grids}(d)--(f) shows the ten-domain cross-domain grids
for the two single-stage encodings and for the exact-rerank ceiling;
Table~\ref{tab:encodings} summarizes them alongside the hardware each
requires. Each grid is read the same way: a row names the domain
whose examples are stored in the bank, a column the domain the
held-out test activations come from -- i.e., a cell scores a bank
built on behavior/domain $X$ against test material from
$Y$.\inote[aiviolet]{Sep 2, your (E): the how-to-read sentence, kept
to one sentence since the cross-domain observation below already
explains the off-diagonal reading (your note that (E) is partly
covered there)} The top row of the same figure places the published
four-domain baseline and the first-mapping Euclidean grid
(Sec.~\ref{subsec:FirstPipeline}) above them for reference; the two
protocols differ in datasets and splits, so the rows of the figure are
not comparable cell by cell.\inote[aiviolet]{Sep 1, your ``why is our
Fig.~2(d)--(f) richer than Fig.~2 in Truthfulness.pdf'': expected,
not a data slip. Their Fig.~2 trains probes on only seven of their
datasets (five truth types + sycophancy + expectation-inverted) and
tests on ten -- roleplaying, insider trading, and sandbagging are
test-only columns for them, since those on-policy benchmarks have no
probe-training split. Our banks are stored labeled exemplars, so a
bank can be built from \emph{any} labeled dataset, and our grid
legitimately has bank rows theirs lacks. Still queued for I.~Bera:
confirm the combined data feeding (d)--(f) is the current run, and
the evidential/ethical naming flag in
Sec.~V-A}\inote[aiviolet]{Sep 2, your ``do we need to repeat/discuss
both rows -- let's game this out'': no figure change made this pass.
The tradeoff: (a) repeats Ying et al.\ and (c) is derivable from
(a)$-$(b), so both are cuttable; (b) is our own Euclidean pipeline
run on their protocol, with no analog in their paper. A middle path
is a top row of just (b). Parked for your call}

\ib{Item 6: these four-domain numbers are not from my runs. The
reference values (1.00/1.00/1.00/0.95) are the Truthfulness Spectrum
paper's own in-domain results; the measured values and the gaps
($-0.03$ here, $-0.06$ in the Aug.~14 document) come from the
undergrads' first pipeline, so I cannot say which of their runs is
current. All my runs are on the ten-domain grid.}
The top row also scores the first mapping on its own terms. Against
the published four-domain baseline, in-distribution \auroc was
1.00/0.99/1.00/0.95 versus the paper's 1.00/1.00/1.00/0.95, the
average cell-wise gap over the full grid was $-0.03$, and the single
worst cell lost 0.09 (combined training, tested on empirical) --
panels (a)--(c) of Fig.~\ref{fig:grids}.\footnote{Those values are
in-distribution \auroc on a four-domain grid and a different protocol
from the ten-domain grid below; Fig.~\ref{fig:grids} keeps the two
protocols in separate rows, and they are not comparable cell by
cell.}\inote[aiviolet]{moved here from old Sec.~IV-A per your
restructure. Standing question from Aug 28: this text matches the
uGrads' deck (mean gap $-0.03$, worst $-0.09$); the Aug 14 methods
document said $-0.06$ / within 0.05 -- worth asking which run is
current}

We make three observations on the ten-domain grids.

\textbf{In-domain fidelity is essentially unaffected by the encoding:}
When a bank is tested on the domain it was built from, the compressed
codes give up almost nothing: the diagonal means are 0.920 (\signrp
ternary, 384~b) and 0.919 (thermometer, 384~b), against 0.928 for the
exact rerank. Said differently, replacing full-precision cosine
similarity over the raw 8{,}192 values with a 384-bit code in a plain
array costs 0.008--0.009 \auroc. The factual block (definitional,
empirical, fictional -- all nine cells, transfer included) averages
0.974 for ternary and 0.942 for thermometer, against the 0.978
ceiling.\inote[aiviolet]{Sep 2: rewritten in plain English per your
note, voice rules applied; every number unchanged}

\textbf{Cross-domain fidelity is where compression shows, and it
separates the two encodings:} An off-diagonal cell scores a bank
built from one domain against test examples from another, so it
measures how well a stored library transfers to material it was not
built from. That transfer is where the encodings pay: off-diagonal
means are 0.640 for ternary and 0.579 for thermometer, against 0.705
for the full-precision ceiling -- a cost of 0.065--0.126 \auroc,
where in-domain the same compression cost less than 0.01. The
ordering is the surprise. At this matched 384-bit point, ternary is
ahead of thermometer by 0.061 \auroc on average and in 63 of 90
cells, which is the opposite of what a preference for the cheaper
cell would predict. The comparison is matched in \emph{bits} but not
in \emph{silicon}, however, since a ternary position may cost roughly
two bits of hardware. At matched cell count -- the thermometer code at
768~b, 256 directions at 3 bits each, same splits -- thermometer scores
0.924 in-domain and 0.631 cross-domain: within 0.009 of ternary's
cross-domain mean, with ternary still ahead in 54 of 90 cells (was 63
of 90 at 384~b) and the worst off-diagonal cell improving from 0.259
to 0.319. At matched silicon the two encodings therefore appear to be
within noise on the mean, with ternary keeping a small cell-wise
edge.\ib{Item 1: matched-cell-count run, thermometer at 768~b (256
directions $\times$ 3 bits), same splits as the 384-b runs. Files:
bits\_grid\_768b\_L4\_m256.csv, matched\_cell\_count.py. This replaces
the ``that experiment has not been run'' caveat and the orange flag
that followed it.}

\textbf{The pattern of failure reproduces the published one:} All three
grids keep the factual block high and put sycophancy near chance from
every other row, which is the qualitative result the source study
reports~\cite{ying2026spectrum}. Sycophancy's own in-domain cell is the
weakest diagonal entry in every configuration (0.627, 0.633, 0.636),
stable across conditions, so it appears to be a property of the data
rather than a compression artifact. The weakest off-diagonal cell sits
at a similar depth in every configuration (0.259, 0.226, 0.220) but
its \emph{location} moves with the encoding:
\texttt{repe\_honesty}~$\to$~sandbagging for ternary,
logical~$\to$~insider trading for thermometer, and
\texttt{repe\_honesty}~$\to$~fictional for the rerank
configurations. The bigger picture is that the accuracy ceiling on
transfer is set by the representations, not by our compression -- the
full-precision grid has the same holes, in the same places. Pooling
does not change this: a bank built from all ten domains' training
splits (11{,}278 rows) averages 0.906 (thermometer, 384~b), 0.909
(ternary, 384~b), and 0.927 (ceiling) across the ten test domains, so
pooling costs nothing against the per-domain banks, but sycophancy
stays at 0.63 -- unlike the source study's combined probe, which
reaches 0.97 there.\ib{Item 2: pooled bank, AUROC per test domain in
the order defn, empir, fict, logic, ethical, syco, repe, roleplay,
insider, sandbag -- thermometer-384: 0.936 0.976 0.966 0.803 0.977
0.627 0.896 0.887 0.997 0.998 (mean 0.906); ternary-384: 0.937 0.976
0.964 0.747 0.977 0.635 0.936 0.919 0.997 1.000 (mean 0.909); cosine
ceiling: 0.970 0.996 0.983 0.803 0.983 0.636 0.957 0.941 0.997 1.000
(mean 0.927). File: combined\_bank.csv (combined\_bank\_auroc.py; it
also has a leave-one-domain-out row).} Covering
a behavior suite at this accuracy therefore means storing directions
for each behavior and domain rather than hoping one direction
travels, i.e., operating at the library scale this paper
prices.\inote[aiviolet]{Aug 28: sanity check against the
plotted arrays, per your student's flag. The earlier text claimed
repe$\to$fictional was the worst cell in every configuration; the
\emph{values} 0.259/0.226/0.220 are each configuration's minimum, but
the cell differs for ternary and thermometer. Every other quoted
statistic in this section verified against the arrays to the stated
precision. CLOSED Aug 28 night against his source files: all four
plotted grids match the methods document cell-for-cell (zero
differences in all 400 cells), and the transposed comparison fails in
all 90 off-diagonal cells per grid, so there is no train/test
transposition -- the figure plots his data exactly. His document's own
min/max/mean lines also agree with every aggregate quoted here, and it
names repe$\to$fictional as the min driver only for the two rerank
grids, consistent with the corrected text above}

\inote[aiviolet]{Aug 28: the top row is the ``all of this in 1
figure'' request. Values transcribed cell-by-cell from slide 4 of the
uGrads' deck (Gisiner/Haidinger/Whalen); every cell was present, so no
white placeholder squares were needed. The $\Delta$ panel recomputed
from (a) and (b) matches the deck's printed $\Delta$ in all 20 cells,
a consistency check on the transcription. Panel (b) is the same 3-bit
LSQ Euclidean pipeline Sec.~IV-B describes -- its diagonal
1.00/0.99/1.00/0.95 matches the values quoted there. Cut (c) if the
figure needs to shrink; (a) and (b) carry the claim}

\begin{table}[t]
\centering
\caption{The encodings, their measured fidelity on the ten-domain grid,
and what each needs from a cell. Bank size is for a 4{,}000-row exemplar
bank at the given code width. The in-domain minimum is sycophancy in
every row; the cross-domain minimum is 0.259 / 0.226 / 0.319 / 0.219 /
0.220 top to bottom, and its cell location varies by encoding
(Sec.~\ref{subsec:Grids}). $^{*}$Matched cell count against ternary at
384~b.}
\label{tab:encodings}
{\scriptsize\setlength{\tabcolsep}{3.5pt}
\begin{tabular}{@{}lccll@{}}
\toprule
\textbf{Encoding} & \textbf{in-dom.} & \textbf{cross-dom.} &
\textbf{Cell needed} & \textbf{Bank} \\
 & mean & mean & & (4K rows) \\
\midrule
\signrp ternary, 384\,b       & 0.920 & 0.640 & \tcam, $\sim$16T   & 192\,kB \\
Thermometer, 384\,b           & 0.919 & 0.579 & BCAM, $\sim$6--10T & 192\,kB \\
Thermometer, 768\,b$^{*}$      & 0.924 & 0.631 & BCAM, $\sim$6--10T & 384\,kB \\
Shortlist $+$ rerank          & 0.928 & 0.671 & BCAM $+$ FP unit   & 192\,kB $+$ raw \\
Exact rerank (ceiling)        & 0.928 & 0.705 & FP baseline        & 65\,MB \\
\bottomrule
\end{tabular}}
\end{table}

\subsection{Cross-domain transfer is hard, and wider codes do not
close the gap}
\label{subsec:BitWidth}
\begin{figure}[!t]
\centering
\includegraphics[width=3.4in]{fig_bitwidth_v1.pdf}
\caption{The cost of a wider code, cell by cell: the change in
cross-domain \auroc as the \signrp ternary code grows from 384 to
2{,}048 bits (red is fidelity lost; off-diagonal cells losing more
than 0.05 are outlined). The diagonal is flat to slightly positive
while the transfer losses concentrate in the rows already weakest at
384 bits, so the extra capacity sharpens domain-specific boundaries
rather than helping them travel (Sec.~\ref{subsec:BitWidth}).}
\label{fig:bitwidth}
\end{figure}

\begin{figure*}[t]
\centering
\includegraphics[width=7.0in]{fig_shortlist_v1.pdf}
\caption{Why a fixed shortlist size is a hardware problem.
(a) Coverage: how often the true full-precision top-1 neighbor is inside
the $k=400$ Hamming shortlist, for two banks that differ by
$\sim$4$\times$ in size, with each bank's chance floor $k/N$ marked.
(b) Cross-domain mean \auroc as $k$ grows as a fraction of each domain's
own bank. Recovery is real but costs most of the bank.}
\label{fig:shortlist}
\end{figure*}
The transfer gap of Sec.~\ref{subsec:Grids} raises the natural
question of whether more code capacity would close it -- and the
answer matters, because if it does not, covering more behaviors means
storing \emph{more} comparison vectors rather than wider ones, i.e.,
growing $K$. Widening the code should tighten the angle estimate in
Eq.~(\ref{eq:signrp}) and improve matching, and in-domain it does.
Cross-domain it does not, in either encoding.\inote[aiviolet]{Aug 31:
reframed per your bottom-margin note (``cross-domain transfer is hard
so we need more comparison vectors''). No push-back on the substance
-- it is the same argument the red-team appendix already makes -- but
I kept the claim conditional on this one model and grid, as objection
(1) requires} Growing the code from 384
to roughly 2{,}048 bits, 23 of 90 off-diagonal cells drop by more than
0.05 \auroc for the ternary code, with a worst single-cell change of
$-0.303$ (sandbagging tested on insider trading); for the thermometer
code the same growth moves 12 of 90 cells by more than 0.05, worst
$-0.209$. Fig.~\ref{fig:bitwidth} maps the change cell by cell for the
ternary code: the diagonal is flat or slightly up, the losses
concentrate in the rows that transfer worst at 384 bits
(\texttt{repe\_honesty}, sandbagging, sycophancy), and the one column
that improves broadly -- insider trading -- is a domain, not the
encoding, getting easier to hit.

A plausible reading is that extra dimensions preserve more of the
geometry of the training domain, including the parts that do not
generalize, so a wider code is a more faithful encoding of a
domain-specific decision boundary. Whatever the mechanism, the
takeaway for hardware is simple: {\bf (1)} the cheaper array is also
the one that transfers better -- on this grid, the costlier cell does
not buy accuracy -- so there is no argument for
paying more silicon per row; and {\bf (2)} if transfer must come from
somewhere, it comes from storing more vectors, not wider ones --
which is the library growth Sec.~\ref{sec:Scaling} prices.

The same holds for a richer cell once the projection is frozen.
Keeping everything else fixed -- the same frozen random projection to
128 values, the same splits and $k=31$ -- and rounding each value to
eight levels (3 bits), a Euclidean matchline (128 multi-bit cells per
row) scores 0.920 in-domain and 0.579 cross-domain, against 0.920 and
0.593 for Hamming distance on the corresponding thermometer code (896
binary cells per row), and 0.916 and 0.598 for Euclidean on the
unrounded values; Euclidean is ahead in only 37 of 90 cells, and at
four levels the two are equal ($-0.003$). Rounding to 3 bits costs
about 0.02. The accuracy edge of the first pipeline
(Sec.~\ref{subsec:FirstPipeline}) therefore came from its trained
projection rather than from its Euclidean matchline; with a frozen
projection, the multi-bit cell's remaining benefit is fewer cells per
row.\ib{Item 8(a), the recombination experiment: same splits, same
frozen random projection to 128 values, same $k=31$ vote, only the
distance changed. Each of the 128 values rounded to one of 8 levels
(3 bits); scored three ways -- Hamming (binary CAM; 896 cells per
row), Euclidean (multi-bit cell; 128 cells per row), and Euclidean on
the unrounded values as a reference.}

\inote[aiviolet]{Aug 28 night, from the student's files: the 2{,}048-bit
\emph{ternary} grid exists (slide 2 of his older deck; saved as
ternary2048.npy), so this $\Delta$ figure per your pick; the worst cell
is sandbagging$\to$insider at $-0.303$, now named in the text. The
2{,}048-bit \emph{thermometer} grid is still aggregate-only (12/90,
$-0.209$). One count note: my recount from his 3-decimal grids gives
24 cells past $-0.05$, not 23 -- two cells sit at exactly $-0.051$
(ethics$\to$evidential, syco$\to$roleplay) -- so his 23 presumably
comes from unrounded data; text keeps his count, and the figure
outlines 24}

\subsection{A fixed shortlist does not travel across bank sizes}
\label{subsec:Shortlist}
The two-stage variant of Sec.~\ref{subsec:TwoStage} is the natural way
to recover the ceiling, and at first reading it works: with an absolute
shortlist of $k=400$, the diagonal mean is unchanged at 0.928 and the
off-diagonal mean falls only from 0.705 to 0.671. The aggregate hides a
failure of a \emph{fixed} parameter, which matters more for a hardware
design than for a software one.

Bank sizes differ by roughly $4\times$ across domains, from 831 rows
(definitional) to 3{,}201 (sycophancy). Fig.~\ref{fig:shortlist}(a)
measures coverage: how often the true full-precision top-1 neighbor
actually lands inside the Hamming shortlist. For the definitional bank
the mean is 82.5\%, against a chance floor of $k/N = 48.1\%$. For the
sycophancy bank the mean is 18.3\% against a floor of 12.5\%, i.e.,
barely above chance, and several test domains sit at or below their own
floor. Consequently, 8 of the 10 training rows are stable between $k=400$
and full precision while the sycophancy row loses 0.204 \auroc on
average and up to 0.468 in a single cell; the ethical row shows a smaller
version of the same effect at $-0.079$.

Fig.~\ref{fig:shortlist}(b) shows this is a sizing problem rather than a
broken mechanism. Sweeping $k$ as a fraction of each domain's own bank,
sycophancy's cross-domain mean recovers from 0.532 at 12.5\% to 0.736 at
100\%, mostly monotonically, and the 100\% value matches the
independently computed full-precision row average to three decimal places
(0.736 vs.\ 0.736; definitional 0.750 vs.\ 0.751), which is a useful
self-consistency check on the pipeline. The recovery costs 75--100\% of
the bank, at which point it is not a shortlist. For a \cam-resident
design, that argues for the single-stage encodings above, or for a
threshold rather than a top-$k$ criterion so the number of survivors
adapts to the bank instead of being fixed at design time.

\subsection{The destroyed-vector audit}
\label{subsec:Audit}
\label{subsec:AuditResults}
\begin{table}[t]
\centering
\caption{The destroyed-vector audit across five variants of the first
pipeline (July 30). Front-end share is the accuracy still present after
every stored row is replaced by a random vector, i.e., the share the
trained layers were carrying alone. High is bad: 100\% means the memory
contributed nothing. This audit is our sanity check that a pipeline is
not ``cheating,'' i.e., classifying in a trained front end rather than
in the stored bank; we run it on every configuration reported in this
paper.}
\label{tab:audit}
{\footnotesize
\begin{tabular}{@{}lcl@{}}
\toprule
\textbf{Variant} & \textbf{Front-end share} & \textbf{Outcome} \\
\midrule
Learned binary hash          & 98--100\% & retired \\
Weighted thermometer         & 66--85\%  & retired \\
3-bit LSQ, Euclidean         & 24--28\%  & survives \\
Direct sign, no projection   & 20\%      & survives \\
Random sign projection       & 4--14\%   & survives \\
\bottomrule
\end{tabular}}
\end{table}
Any pipeline of the form ``trained front end, then memory'' can pass
its accuracy test while the memory contributes nothing. We therefore
re-measure every configuration with the stored rows replaced by random
vectors of the same shape: a \cam-based detector should collapse to
chance when the bank is destroyed, and if it does not, whatever sits
in front of the array was doing the classification. We report the
accuracy that survives destruction as the \emph{front-end share},
where high is bad: 100\% means the memory was decoration.

The audit has changed decisions. Table~\ref{tab:audit} reports it
across five variants of the first pipeline
(Sec.~\ref{subsec:FirstPipeline}). In the two highest-scoring
variants, one of which reached 0.993 \auroc, between 66\% and 100\%
of the accuracy survived randomizing every stored row: the trained
layers in front of the array were performing most of the
classification, and the \cam was close to decoration. Both were
retired. The 3-bit LSQ Euclidean configuration sits at 24--28\%, so
its memory is contributing; the variants whose stored rows carry
nearly everything are the ones with a frozen or absent projection, at
4--20\%, which is why the encodings of Secs.~\ref{subsec:SignRP}
and~\ref{subsec:Thermometer} use frozen projections -- with nothing
trained in front of the array, the failure has nowhere to occur.

The encodings reported above pass. Randomizing every stored
activation collapses the 2{,}048-b \signrp grid to a mean of 0.50,
with 83\% of cells between 0.4 and 0.6, and a label-shuffle control
at 384~b gives 0.500 for ternary (all 100 cells within 0.40--0.56) and
0.501 for thermometer (all cells within 0.44--0.57), so the
classification is coming from what is stored rather than from
anything in front of the array.\ib{Item 5: the 0.50 / 83\% figures are
from the 2{,}048-b run, now said so. At 384~b I have the label-shuffle
audit -- mean 0.500 ternary (all cells 0.40--0.56), 0.501 thermometer
(0.44--0.57) -- which is a different control from the randomized-bank
audit; the text states the two separately. If a randomized-bank run at
384~b is wanted, it is a short rerun.} An
accuracy number alone would not have separated the five variants of
Table~\ref{tab:audit} -- the highest-scoring one was the worst
offender. As such, we now run the audit by default on every
configuration, report front-end share alongside \auroc, and treat a
variant whose accuracy survives bank destruction as an accelerator
result rather than a memory result. The recombined design of
Sec.~\ref{subsec:PathBack}, now run (Sec.~\ref{subsec:BitWidth}),
passes by construction -- nothing is trained ahead of its array.

\inote[aiviolet]{Aug 28 night: verified against the full audit grid in
the student's older deck (slide 3) -- mean 0.504, 83.0\% of cells in
[0.4, 0.6], so the quoted figures are exact. That grid is the
2{,}048-bit configuration, though, so the standing question is
narrower but still open: confirm the audit was also run at 384 bits,
or say ``at 2{,}048 bits'' in this sentence}
